\chapter{Skema Penyimpanan Data RME}
\label{lampiran:skema-penyimpanan-rme}

Lampiran ini memuat rincian skema penyimpanan data RME yang digunakan pada rancangan segmentasi data. Tabel~\ref{tab:onchain-schema} dan Tabel~\ref{tab:offchain-schema} masing-masing menjabarkan metadata segmen pada penyimpanan \textit{on-chain} dan isi segmen terenkripsi pada penyimpanan \textit{off-chain}.

\begin{longtable}{|>{\raggedright\arraybackslash}p{2.8cm}|>{\raggedright\arraybackslash}p{4.0cm}|>{\raggedright\arraybackslash}p{5.4cm}|}
	\caption{Skema Metadata Segmen RME \textit{On-Chain}} \label{tab:onchain-schema}                                                              \\
	\hline
	\textbf{Komponen} & \textbf{Key JSON}                                                                                   & \textbf{Keterangan} \\
	\hline
	\endfirsthead
	\hline
	\textbf{Komponen} & \textbf{Key JSON}                                                                                   & \textbf{Keterangan} \\
	\hline
	\endhead

	\multirow{4}{2.8cm}{\textbf{Identitas Segmen}}
	                  & \texttt{segment\_id}
	                  & UUID unik untuk satu segmen data RME.                                                                                     \\ \cline{2-3}
	                  & \texttt{related\_rme\_id}
	                  & ID RME atau episode pelayanan yang menjadi induk segmen, dengan format \verb|RME-{6_digit}|.                              \\ \cline{2-3}
	                  & \texttt{patient\_address}
	                  & Alamat \textit{wallet} pasien pemilik RME.                                                                                \\ \cline{2-3}
	                  & \texttt{hospital\_id}
	                  & ID fasilitas pelayanan kesehatan ketika menerima pelayanan.                                                               \\ \hline

	\multirow{2}{2.8cm}{\textbf{Kategori Data}}
	                  & \texttt{dataset\_category}
	                  & Kategori dataset. Sesuai dengan Tabel~\ref{tab:mapping-dataset-category}.                                                 \\ \cline{2-3}
	                  & \texttt{function\_category}
	                  & Kategori fungsi klinis atau administratif. Sesuai dengan Tabel~\ref{tab:mapping-function-category}.                       \\ \hline

	\multirow{2}{2.8cm}{\textbf{Pointer dan Integritas}}
	                  & \texttt{ipfs\_cid}
	                  & CID IPFS dari segmen data terenkripsi.                                                                                    \\ \cline{2-3}
	                  & \texttt{integrity\_hash}
	                  & Hash SHA-256 dari segmen terenkripsi.                                                                                     \\ \hline

	\multirow{2}{2.8cm}{\textbf{Kapsul Kunci}}
	                  & \texttt{capsule}
	                  & Kapsul PRE yang diperlukan untuk proses \textit{re-encryption}.                                                           \\ \cline{2-3}
	                  & \texttt{enc\_key\_and\_nonce}
	                  & Kunci simetris dan \textit{nonce} yang sudah terenkripsi untuk membuka segmen data.                                       \\ \hline

	\multirow{2}{2.8cm}{\textbf{Audit Metadata}}
	                  & \texttt{created\_at}
	                  & Waktu pembuatan metadata segmen.                                                                                          \\ \cline{2-3}
	                  & \texttt{author\_address}
	                  & Alamat \textit{wallet} aktor yang membuat atau memperbarui segmen.                                                        \\ \hline
\end{longtable}

\begin{longtable}{|>{\raggedright\arraybackslash}p{2.8cm}|>{\raggedright\arraybackslash}p{4.0cm}|>{\raggedright\arraybackslash}p{5.4cm}|}
	\caption{Skema Segmen Data RME \textit{Off-Chain}} \label{tab:offchain-schema}                                                                                     \\
	\hline
	\textbf{Komponen} & \textbf{Key JSON}                                                                                                        & \textbf{Keterangan} \\
	\hline
	\endfirsthead
	\hline
	\textbf{Komponen} & \textbf{Key JSON}                                                                                                        & \textbf{Keterangan} \\
	\hline
	\endhead

	\multirow{5}{2.8cm}{\textbf{Header Segmen}}
	                  & \texttt{segment\_id}
	                  & UUID yang sama dengan \texttt{segment\_id} pada metadata \textit{on-chain}.                                                                    \\ \cline{2-3}
	                  & \texttt{related\_rme\_id}
	                  & ID RME atau episode pelayanan yang menjadi induk segmen, dengan format \verb|RME-{6_digit}|.                                                   \\ \cline{2-3}
	                  & \texttt{dataset\_category}
	                  & Kategori dataset sesuai Tabel~\ref{tab:mapping-dataset-category}.                                                                              \\ \cline{2-3}
	                  & \texttt{function\_category}
	                  & Kategori fungsi sesuai dengan Tabel~\ref{tab:mapping-function-category}.                                                                       \\ \hline

	\multirow{4}{2.8cm}{\textbf{Konteks Medis}}
	                  & \texttt{patient\_address}
	                  & Alamat IOTA pasien pemilik RME.                                                                                                                \\ \cline{2-3}
	                  & \texttt{service\_date}
	                  & Tanggal dan waktu ketika segmen tercatat dengan format \texttt{YYYY-MM-DD HH:MM:SS} beserta informasi \textit{timezone}.                       \\ \cline{2-3}
	                  & \texttt{author\_address}
	                  & Alamat IOTA aktor yang menulis isi segmen.                                                                                                     \\ \hline

	\multirow{2}{2.8cm}{\textbf{Isi Segmen}}
	                  & \texttt{payload}
	                  & Objek segmen data RME dengan tipe data \texttt{Value}. Pada implementasi ini, isi segmen dalam bentuk \texttt{text}.                           \\ \cline{2-3}
	                  & \texttt{payload\_hash}
	                  & \textit{Hash} isi \texttt{payload}.                                                                                                            \\ \hline
\end{longtable}